% PDF p. 270
\BackSection{习题}
\begin{enumerate}
  \item \textbf{习题 5.1}\par 写出一阶偏导：
  \[
    \begin{aligned}
      X_u&:=\frac{\partial X}{\partial u}
            =\left(\frac{\partial x}{\partial u},\frac{\partial y}{\partial u},
                    \frac{\partial z}{\partial u}\right),\\
      X_v&:=\frac{\partial X}{\partial v}
            =\left(\frac{\partial x}{\partial v},\frac{\partial y}{\partial v},
                    \frac{\partial z}{\partial v}\right).
    \end{aligned}
  \]
  计算 $EG-F^2$。
  \item \textbf{习题 5.2}\par 令 $z=\frac12(u^2-v^2)$。计算
  \[
    e=\langle X_{uu},N\rangle,\qquad
    f=\langle X_{uv},N\rangle,\qquad g=\langle X_{vv},N\rangle.
  \]
  \item \textbf{习题 5.3}\par 请证明以下陈述：“把关于 $s$ 的导数解释为沿变分的协变导数（再在 $s=0$ 处求值），得到
  \[
    \left.\nabla_{\partial_s}\nabla_{\partial_t}\partial_t\gamma_s(t)\right|_{s=0}=0.\text{”}
  \]
  \item \textbf{习题 5.4}\par 计算三维球面的里奇张量 $R_{ij}$。
  \item \textbf{习题 5.5}\par 证明对称性：由无挠性得 $\mathcal A(X,Y)=\mathcal A(Y,X)$。
  \item \textbf{习题 5.6}\par 计算球面 $x^2+y^2+z^2=R^2$ 的主曲率和高斯曲率。
  \item \textbf{习题 5.7}\par 计算鞍形曲面 $z=x^2-y^2$ 在原点处的主曲率和高斯曲率。
  \item \textbf{习题 5.8}\par 计算椭球面 $\frac{x^2}{a^2}+\frac{y^2}{b^2}+\frac{z^2}{c^2}=1$ 的主曲率和高斯曲率。
  \item \textbf{习题 5.9}\par 计算半径为 $r=2$ 的四维球面的里奇张量和标量曲率。
  \item \textbf{习题 5.10}\par 计算庞加莱圆盘模型在 $r=0.9$ 处的里奇张量。
\end{enumerate}

% PDF p. 271
\needspace{12\baselineskip}
\BackSection{参考文献}
\begin{bookreferences}
  \item Hirvonen, V. (2025). The Ricci Tensor: A complete Guide with Examples.
    \url{https://profoundphysics.com/the-ricci-tensor/}
  \item Jia, Y.-B. (2024). Geodesics (COMS 4770/5770 Notes).
    \url{https://faculty.sites.iastate.edu/jia/files/inline-files/geodesics.pdf}
  \item Koch, R. (2005). Lectures on Gauss-Bonnet. May 30, 2005.
    \url{https://pages.uoregon.edu/koch/math433/Gauss_Bonnet.pdf}
  \item Levi-Civita, T. (1926). The absolute differential calculus (calculus of tensors). Blackie \& Son.
\end{bookreferences}
